% ============================================================
%  Mesa vs Modelly — Présentation Beamer (version corrigée)
%  Compilez avec : xelatex mesa_vs_modelly.tex  (2 passes)
% ============================================================
\documentclass[aspectratio=169, 11pt]{beamer}

% ---------- thème ------------------------------------------
\usetheme{Madrid}
\usecolortheme{seahorse}
\usefonttheme{professionalfonts}
\setbeamertemplate{navigation symbols}{}

% ---------- encodage & langue ------------------------------
\usepackage{fontspec}
\usepackage[french]{babel}
\usepackage{csquotes}
\usepackage{booktabs}

% ---------- code -------------------------------------------
\usepackage{listings}
\usepackage{xcolor}

\definecolor{codebg}{RGB}{246,246,250}
\definecolor{keyword}{RGB}{30,100,200}
\definecolor{comment}{RGB}{100,120,100}
\definecolor{string}{RGB}{180,40,40}
\definecolor{good}{RGB}{20,140,60}
\definecolor{bad}{RGB}{190,40,30}
\definecolor{warncolor}{RGB}{200,120,0}

\lstset{
  backgroundcolor=\color{codebg},
  basicstyle=\ttfamily\small,
  keywordstyle=\color{keyword}\bfseries,
  commentstyle=\color{comment}\itshape,
  stringstyle=\color{string},
  language=Python,
  breaklines=true,
  frame=none,
  showstringspaces=false,
  tabsize=2,
  columns=flexible,
  xleftmargin=4pt,
  xrightmargin=4pt,
  morekeywords={self,True,False,None,super},
}

% ---------- icônes -----------------------------------------
\newcommand{\okmark}{\textcolor{good}{\textbf{\checkmark}}}
\newcommand{\komark}{\textcolor{bad}{\textbf{$\boldsymbol{\times}$}}}
\newcommand{\warnmark}{\textcolor{warncolor}{\textbf{!}}}

% ---------- métadonnées ------------------------------------
\title[Mesa \textrightarrow{} Modelly]{Pourquoi nous avons migré de Mesa vers Modelly}
\subtitle{Modélisation individu-centrée — projet TIS}
\author{Marine Courtois}
\institute{Modelly-TIS-Research}
\date{Novembre 2026}

% ============================================================
\begin{document}

% ----------------------------------------------------------
\begin{frame}
  \titlepage
\end{frame}

\begin{frame}{Plan}
  \tableofcontents
\end{frame}

% ============================================================
\section{Rappel — modèles individu-centrés}

\begin{frame}{C'est quoi un modèle individu-centré ?}

  \begin{block}{Principe}
    Chaque individu est programmé avec ses propres règles de comportement.
    Le comportement \textbf{global} du système émerge des interactions locales.
  \end{block}

  \bigskip

  \begin{columns}[T]
    \column{0.46\textwidth}
      \textbf{Nos agents dans le modèle TIS}
      \medskip
      \begin{itemize}\setlength\itemsep{4pt}
        \item Mâles sauvages (\textit{D. suzukii})
        \item Mâles stériles lâchés (TIS)
        \item Femelles — accouplement, ponte
        \item Larves
        \item Fruits — statut, ovipositions
      \end{itemize}

    \column{0.50\textwidth}
      \textbf{Ce que simule notre modèle}
      \medskip
      \begin{itemize}\setlength\itemsep{4pt}
        \item Une saison agricole jour par jour
        \item Compétition stériles\,/\,sauvages
        \item Dommages larvaires selon la stratégie de lâcher
        \item 2 cultures : fraise et cerise
      \end{itemize}
  \end{columns}

\end{frame}

% ============================================================
\section{Présentation des deux packages}

\begin{frame}{Mesa — le package historique}

  \textbf{Créé en 2015} — le framework ABM de référence en Python

  \bigskip

  \begin{columns}[T]
    \column{0.60\textwidth}
      \begin{itemize}\setlength\itemsep{6pt}
        \item Open-source, licence Apache 2.0
        \item Très utilisé en sciences sociales et écologie
        \item Grande communauté, nombreux tutoriels
        \item Version actuelle : 3.5.3
      \end{itemize}

    \column{0.36\textwidth}
      \begin{alertblock}{Point de vigilance}
        L'API a subi des \textbf{ruptures majeures} entre la v2 et la v3 qui cassent les modèles existants.
      \end{alertblock}
  \end{columns}

\end{frame}

\begin{frame}{Modelly — orienté reproductibilité scientifique}

  \textbf{Package ABM plus récent} — pensé pour les simulations traçables

  \bigskip

  \begin{columns}[T]
    \column{0.60\textwidth}
      \begin{itemize}\setlength\itemsep{6pt}
        \item Open-source, licence Apache 2.0
        \item API très proche de Mesa — migration facile
        \item Accent sur reproductibilité et études paramétriques
        \item Version actuelle : 1.5.0
      \end{itemize}

    \column{0.36\textwidth}
      \begin{block}{Même base que Mesa}
        \texttt{numpy}, \texttt{pandas},\\
        \texttt{scipy}, \texttt{tqdm}\\
        \medskip
        Aucune dépendance\\supplémentaire
      \end{block}
  \end{columns}

\end{frame}

% ============================================================
\section{Trois raisons scientifiques de migrer}

% ----------------------------------------------------------
\begin{frame}{Raison \textnumero{}1 — Reproductibilité}

  \begin{alertblock}{Le problème avec Mesa}
    Deux sources aléatoires séparées : le modèle \textbf{et} le scheduler.\\
    La séquence des tirages a changé entre Mesa 2.x et 3.x.
  \end{alertblock}

  \bigskip

  \begin{block}{La solution Modelly : \texttt{FastRandom}}
    Un seul générateur — \texttt{model.random} — Mersenne Twister stable entre versions.\\[6pt]
    Deux runs avec la même graine produisent des résultats \textbf{identiques bit-à-bit}.
  \end{block}

  \bigskip

  \okmark\; Vérifié automatiquement par \texttt{tests/test\_regression.py} à chaque commit.

\end{frame}

\begin{frame}[fragile]{Raison \textnumero{}1 — Reproductibilité (code)}

  \begin{columns}[T]
    \column{0.46\textwidth}
      \textbf{Mesa — graine instable}
      \begin{lstlisting}
# API deprecated en v3
model.reset_randomizer(42)
# FutureWarning !
# Sequence peut changer
# selon la version de Mesa
      \end{lstlisting}

    \column{0.50\textwidth}
      \textbf{Modelly — graine intégrée}
      \begin{lstlisting}
# Graine dans le constructeur
model = RematingModel(seed=42)

# Bit-identique garanti
# entre toutes les versions
      \end{lstlisting}
  \end{columns}

\end{frame}

% ----------------------------------------------------------
\begin{frame}{Raison \textnumero{}2 — Contrôle de l'activation}

  Dans un ABM, \textbf{l'ordre dans lequel les agents agissent change les résultats}.

  \bigskip

  \begin{columns}[T]
    \column{0.46\textwidth}
      \begin{alertblock}{Mesa — boîte noire}
        Le scheduler \texttt{RandomActivation} cache l'ordre d'activation.\\[4pt]
        \komark\; \textbf{Supprimé} dans Mesa 3.0 !
      \end{alertblock}

    \column{0.50\textwidth}
      \begin{block}{Modelly — activation explicite}
        L'ordre est écrit \textbf{directement dans le code} du modèle : lisible, auditable, reproductible.\\[4pt]
        \okmark\; N'importe quel collaborateur peut le lire.
      \end{block}
  \end{columns}

\end{frame}

\begin{frame}[fragile]{Raison \textnumero{}2 — Contrôle de l'activation (code)}

  \begin{columns}[T]
    \column{0.46\textwidth}
      \textbf{Mesa 2.x}
      \begin{lstlisting}
# Scheduler opaque
self.schedule = \
  RandomActivation(self)

def step(self):
    self.schedule.step()
# Supprime en Mesa 3 !
      \end{lstlisting}

    \column{0.50\textwidth}
      \textbf{Modelly — notre advance\_day()}
      \begin{lstlisting}
agents = [a for b in
  self.agents_by_type.values()
  for a in b]
self.random.shuffle(agents)
for a in agents:
    if a in self.agents:
        a.step()
      \end{lstlisting}
  \end{columns}

\end{frame}

% ----------------------------------------------------------
\begin{frame}{Raison \textnumero{}3 — Stabilité de l'API}

  \textbf{Mesa a cassé l'API de façon répétée entre v2 et v3 :}

  \bigskip

  \begin{center}
  \begin{tabular}{lcc}
    \toprule
    Ce qui a changé & Version & Impact \\
    \midrule
    Signature de \texttt{Agent.\_\_init\_\_} & 2 $\to$ 3 & \komark\; Cassant \\
    \texttt{RandomActivation} supprimé & 2 $\to$ 3 & \komark\; Cassant \\
    \texttt{mesa.time.*} remplacé & 2 $\to$ 3 & \komark\; Cassant \\
    \texttt{reset\_randomizer()} déprécié & 3.x & \warnmark\; Warning \\
    \bottomrule
  \end{tabular}
  \end{center}

  \bigskip

  \begin{block}{Modelly}
    API rétrocompatible depuis la v1.0, avec un cycle de dépréciation annoncé.
  \end{block}

\end{frame}

% ============================================================
\section{Impact concret sur notre modèle}

\begin{frame}{Qu'est-ce qui a changé dans le code ?}

  \begin{center}
  \textbf{Uniquement la plomberie du framework — pas les équations}
  \end{center}

  \medskip

  \begin{center}
  \begin{tabular}{lll}
    \toprule
    Élément & Mesa (avant) & Modelly (après) \\
    \midrule
    Import & \texttt{import mesa} & \texttt{import modelly} \\
    Constructeur agent & \texttt{\_\_init\_\_(self, uid, model)} & \texttt{\_\_init\_\_(self, model)} \\
    Supprimer un agent & \texttt{schedule.remove(self)} & \texttt{self.remove()} \\
    Graine RNG & \texttt{reset\_randomizer(42)} & \texttt{Model(rng=42)} \\
    Scheduler & \texttt{RandomActivation} & Boucle explicite \\
    \bottomrule
  \end{tabular}
  \end{center}

\end{frame}

\begin{frame}{Ce qui n'a \textbf{pas} changé}

  \begin{columns}[T]
    \column{0.46\textwidth}
      \begin{block}{\okmark\; Inchangé}
        \begin{itemize}\setlength\itemsep{5pt}
          \item Toutes les \textbf{équations biologiques}
          \item Tous les \textbf{tirages aléatoires}
          \item Mortalité, ponte, accouplement
          \item Calendriers de lâcher
          \item Métriques de fin de saison
        \end{itemize}
      \end{block}

    \column{0.50\textwidth}
      \begin{alertblock}{La preuve : tests de non-régression}
        \texttt{make test} compare les sorties\\du moteur Modelly\\
        aux CSV du \textbf{notebook original}.\\[8pt]
        Résultat : \textcolor{good}{\textbf{égalité bit-à-bit \okmark}}
      \end{alertblock}

      \bigskip
      {\small\itshape « La migration n'est pas une réécriture du modèle. C'est un remplacement de plomberie, vérifié automatiquement. »}
  \end{columns}

\end{frame}

% ============================================================
\section{Performances et fonctionnalités bonus}

\begin{frame}{Modelly est aussi plus rapide}

  \textbf{Benchmark mesuré sur la machine de développement (macOS, Python 3.13) :}

  \bigskip

  \begin{center}
  \begin{tabular}{lccc}
    \toprule
    Configuration & Mesa & Modelly & Gain \\
    \midrule
    500 agents $\times$ 30 jours & 3,0 ms & 1,1 ms & \textcolor{good}{\textbf{$\times$2,7}} \\
    9\,500 agents $\times$ 100 jours & 60 ms & 46 ms & \textcolor{good}{\textbf{$\times$1,3}} \\
    \bottomrule
  \end{tabular}
  \end{center}

  \bigskip

  Sur nos études paramétriques (centaines de réplicats), le gain $\times$1,3 se cumule significativement.

\end{frame}

\begin{frame}{Fonctionnalités supplémentaires dans Modelly}

  \begin{columns}[T]
    \column{0.46\textwidth}
      \textbf{DataCollector enrichi}
      \begin{itemize}\setlength\itemsep{5pt}
        \item Export \texttt{Arrow} / \texttt{Parquet}
        \item Export \texttt{xarray}
        \item \texttt{get\_agenttype\_vars\_dataframe()}
      \end{itemize}

    \column{0.50\textwidth}
      \textbf{Infrastructure}
      \begin{itemize}\setlength\itemsep{5pt}
        \item \texttt{model.save()} / \texttt{model.load()} — checkpoint
        \item \texttt{run\_batch()} — parallélisation native
        \item \texttt{run\_grid()} $\to$ \texttt{DataFrame} directement
      \end{itemize}
  \end{columns}

\end{frame}

% ============================================================
\section{Conclusion}

\begin{frame}{Conclusion — Pourquoi Modelly ?}

  \begin{columns}[T]
    \column{0.30\textwidth}
      \begin{block}{\textnumero{}1 Reproductibilité}
        Un seul RNG (\texttt{FastRandom}), stable entre versions. Tests automatiques.
      \end{block}

    \column{0.30\textwidth}
      \begin{block}{\textnumero{}2 Auditabilité}
        L'ordre d'activation des agents est explicite dans le code.
      \end{block}

    \column{0.30\textwidth}
      \begin{block}{\textnumero{}3 Stabilité}
        API rétrocompatible. Pas de rupture forcée à chaque mise à jour.
      \end{block}
  \end{columns}

  \bigskip\bigskip

  \begin{center}
    \large La migration a changé \textbf{la plomberie}, pas le modèle.\\[6pt]
    \normalsize
    Équations et résultats : \textbf{identiques} \okmark\\
    Reproductibilité et stabilité : \textbf{améliorées} \okmark
  \end{center}

\end{frame}

\begin{frame}{Merci — Questions ?}

  \begin{center}
    \Large\textbf{Merci pour votre attention}
  \end{center}

  \bigskip\bigskip

  \begin{columns}[T]
    \column{0.46\textwidth}
      \textbf{Dans le dépôt}
      \begin{itemize}
        \item \texttt{backend/tis\_engine/model.py}
        \item \texttt{tests/test\_regression.py}
      \end{itemize}

    \column{0.50\textwidth}
      \textbf{Documentation}
      \begin{itemize}
        \item \texttt{remipelloux.github.io}
        \item \texttt{mesa.readthedocs.io}
      \end{itemize}
  \end{columns}

\end{frame}

\end{document}
